3.1 Static Triode Nonlinearity
3.1.1 Physical Motivation and Mathematical Model
The fundamental building block of our vacuum tube simulation is the static nonlinearity that replicates the transfer characteristic of a triode amplifier stage. In a physical triode circuit, the relationship between the grid voltage ($V_g$) and the plate current ($I_p$) is governed by Child-Langmuir law and the specific load line of the circuit. For real-time digital implementation, we adopt a computationally efficient sigmoidal approximation based on the hyperbolic tangent function:

$$ y[n] = \tanh\bigl(a \cdot x[n] + b\bigr) \tag{10} $$

where $x[n]$ is the normalized input signal, $a$ is the drive (or slope) parameter, and $b$ is the bias (or operating point) offset. This model effectively captures the "soft-clipping" behavior of tube saturation, where the gain decreases smoothly as the signal approaches the supply rails (modeled by the $\pm 1$ horizontal asymptotes of the $\tanh$ function).

3.1.2 Parameter Mapping
The two primary parameters, $a$ and $b$, map directly to the tonal characteristics of the stage:

Drive ($a$): Governing the slope of the transfer function in the linear region, $a$ controls the gain and the "hardness" of the saturation. Low values (e.g., $a \approx 1{-}3$) represent clean stages with high headroom, while high values (e.g., $a > 8$) yield early and aggressive saturation typical of high-gain leads.
Bias ($b$): This parameter shifts the operating point of the stage. A zero bias ($b = 0$) results in a point-symmetric transfer function, whereas a non-zero bias ($b \neq 0$) introduces asymmetry. This shift simulates the quiescent grid voltage established by the cathode resistor or external bias supply.

3.1.3 Symmetry and Harmonic Content
The choice of $b$ is critical for the harmonic footprint of the amplifier. As demonstrated in Figure 3, a symmetrical stage ($b=0$) generates predominantly odd-order harmonics (H3, H5, H7...), resulting in a "hollow" or "square-like" distortion characteristic. In contrast, introducing a bias offset ($b \neq 0$) breaks the odd symmetry, generating a rich spectrum of even-order harmonics (H2, H4, H6...). In guitar amplification, the presence of H2 is widely regarded as a primary source of musical "warmth" and "thickness," making the precise calibration of $b$ essential for authentic modeling of vintage circuits.

3.1.4 Real-time Anti-Aliasing (ADAA1)
Nonlinear functions like $\tanh$ generate high-frequency harmonics that can exceed the Nyquist frequency ($f_s/2$), leading to inharmonic aliasing distortion. To mitigate this without the CPU cost of conventional oversampling, we implement the First-order Antiderivative Anti-Aliasing (ADAA1) technique.

In this approach, the $\tanh$ function is replaced by its discrete-time antiderivative average:
$$ y[n] = \frac{f_{ad}(a \cdot x[n] + b) - f_{ad}(a \cdot x[n-1] + b)}{(a \cdot x[n] + b) - (a \cdot x[n-1] + b)} \tag{11} $$
where $f_{ad}(u) = \ln(\cosh(u))$ is the antiderivative of $\tanh(u)$. Mathematically, this acts as a low-pass filter integrated into the nonlinearity itself, significantly suppressing aliased components as shown in Figure 4. For signals with very small changes between samples, the model automatically reverts to the standard $\tanh$ to ensure numerical stability.

3.1.5 Figure Descriptions and Interpretation
Figure 1 — Transfer Function Characteristics. Panel (a) shows how the drive parameter $a$ affects the saturation onset; higher $a$ leads to a more compressed and saturated response for the same input level. Panel (b) illustrates the role of the bias $b$ in breaking the stage's symmetry. By shifting the operating point (red marker), the model creates an imbalance between positive and negative headroom. For a symmetric input signal, this results in unequal distances to the saturation asymptotes, physically modeling the asymmetric clipping that generates even-order harmonics.

Figure 2 --- Waveform Distortion and Contrast Analysis. Figure 2 provides a definitive, high-contrast view of the waveform deformation. Panel (a) employs a "symmetry mirror" technique with a low-drive, high-bias configuration ($a=1.0, b=-1.1$); by overlaying the inverted positive peak over the negative signal, the "asymmetry gap" (purple area) reveals a flagrant disparity in peak shape. One can clearly observe a soft, "breathing" positive peak (rounded) versus a hard, compressed negative peak (flat plateau). Panel (b) illustrates the dynamic compression imbalance: the green shaded area represents the region of soft, progressive rounding where the tube "breathes," while the red area highlights the zone of instant hard saturation where the signal "hits the wall" early.

Figure 3 — Harmonic Content comparison. The frequency-domain analysis validates the symmetry discussion. The symmetrical case (left) shows the expected odd-harmonic dominance. The asymmetrical case (right) reveals the prominent emergence of the second harmonic (H2), which is nearly equal in amplitude to the fundamental and H3, providing the desired "tube character."

Figure 4 — Aliasing Suppression via ADAA1. This figure highlights the benefit of the ADAA1 implementation when processing high-frequency signals. The gray plot (standard $\tanh$) shows numerous "alien" mirrored frequencies (aliasing) caused by high-order harmonics folding back into the audible spectrum. The green plot (ADAA1) shows a much cleaner floor with significant suppression of these inharmonic artifacts, ensuring transparency even at high gain settings and standard sample rates.
